\chapter{Inleiding}

Dit rapport heeft als doel om te onderzoeken op welke manier Jansje op een technisch en financieel haalbare wijze  kan worden verduurzaamd. Dit wordt gedaan in samenwerking met Loods M en andere betrokken stakeholders. De bijdrage van het vervoer over water aan de totale CO$_2$-uitstoot door de Nederlandse economie was bijna 7\% \cite{cbs_water_uitstoot}. Het verduurzamen van schepen zal een belangrijke bijdrage leveren aan het behalen van de klimaatdoelstellingen voor 2050 \cite{imo_ghg_emissions}. \\

\noindent
Loods M is een Fieldlab in Maassluis dat op zoek is naar innovatieve ideeën om schepen te verduurzamen. In Loods M komen studenten, vakmensen en bedrijven samen om slimme ideeën in het echte leven te testen. Eén van de schepen binnen dit project is Jansje. Jansje is een historisch transportschip dat in Maassluis ligt en in 1899 is gebouwd \cite{Jansje}. Het schip is al vijf generaties in bezit van dezelfde familie. Op dit moment heeft Jansje een conventionele dieselmotor, wat valt onder de stage 1b \cite{dieselnet_eu_nonroad_standards}. Deze aandrijving leidt tot de uitstoot van relatief veel emissiegassen. Het verduurzamen van dit schip vormt daarom een relevante technische en maatschappelijke uitdaging. Er worden verschillende aandrijvingen en energiedragers onderzocht om de meest geschikte keuze te maken voor Jansje.\\


\noindent Fieldlab Maassluis heeft samenwerkingen met verschillende bedrijven, echter is niet elk bedrijf relevant in dit onderzoek. Bedrijven die wel relevant zijn, zijn Dr Ten \cite{drten_zeezout_batterij2} die een Zeezoutbatterij zouden kunnen leveren, Nedstack \cite{nedstack_pem_fuel_cell_stacks} die beschikking hebben tot een waterstof brandstofcel. Ook worden er nog andere bedrijven gecontacteerd, voor de nodige informatie. Het onderzoek zal worden gedaan voor het fieldlab en dat zal betekenen dat de voorzieningen van deze bedrijven zullen worden onderzocht maar niet verplicht zijn om te gebruiken.

